% GENERATED FILE. Edit canonical lessons, then run npm run generate:books.
% canonical-source-hash: fnv1a64:5f332d7789aac74c
% canonical-lessons: PA-C201-identity-return, PA-C201-class, PA-C201-class-review, PA-C201-class-label, PA-C201-payoff

\chapter{This Is a Class}
\label{ch:pa-a2-class}
\hlchaptermodality{201}

\begin{chapteropening}
\textbf{By the end of this chapter:} \emph{I can recognise and use the school noun for a class or grade in one bounded identity frame.}

Recognise and produce the school-class noun in one bounded identity frame through listening, speaking, reading, and writing.
\end{chapteropening}

\section[\texorpdfstring{\pa{ਇਹ} · \pa{ਹੈ} (eh; hai)}{ਇਹ · ਹੈ (eh; hai)}]{Return to this and is}

\label{lesson:PA-C201-identity-return}



\begin{quote}
\emph{Hear:} \emph{eh; hai}; YOU SAY: \textbf{\pa{ਇਹ}; \pa{ਹੈ}}
\end{quote}

\begin{grammarlens}[title={Grammar Lens}]
Keep the familiar pointer first and the familiar verb last. The middle slot stays empty until the next lesson owns its new school word.
\end{grammarlens}

\subsection*{Your turn}
\begin{itemize}
\raggedright
  \item \emph{Hear:} \emph{eh}; YOU SAY: \textbf{\pa{ਇਹ}}
  \item \emph{Hear:} \emph{hai}; YOU SAY: \textbf{\pa{ਹੈ}}
  \item Say the first word, pause for the empty middle, then say the last.
\end{itemize}

\begin{tcolorbox}[breakable,colback=teal!4,colframe=teal!35!black,title={Before you move on}]
Produce the two familiar words around the empty middle slot. This opener owns no new word.
\end{tcolorbox}
\section[\texorpdfstring{\pa{ਕਲਾਸ} (kalaas)}{ਕਲਾਸ (kalaas)}]{One school noun for a class or grade}

\label{lesson:PA-C201-class}



\begin{quote}
Retrieve the familiar first and last words of the identity frame.
\end{quote}

\begin{grammarlens}[title={Grammar Lens: direct owner}]
\begin{quote}
\textbf{\pa{ਕਲਾਸ}} — \emph{kalās} — school class; grade
\end{quote}

Hear and say the whole word before reading it. Keep it as one whole school noun rather than inventing smaller Punjabi pieces. This lesson owns the word alone; the complete identity frame comes only after a separate retrieval.
\end{grammarlens}

\begin{cousinweb}
Punjabi University's Punjabipedia explicitly records English \textbf{Class} as the source of modern written \textbf{\pa{ਕਲਾਸ}}. It also lists several meanings; this lesson keeps only the bounded school class-or-grade meaning. The Punjabi school noun comes from English \textbf{Class}. Source: \href{https://punjabipedia.org/topic.aspx?txt=\%E0\%A8\%95\%E0\%A8\%B2\%E0\%A8\%BE\%E0\%A8\%B8}{Punjabi University Punjabipedia}.

The Punjab School Education Board's current circular page uses the word in school scheduling prose. That current school use supports the narrow meaning here. The deeper history and other meanings stay outside this lesson. Source: \href{https://pseb.ac.in/public/pressnotescirculars}{Punjab School Education Board}.
\end{cousinweb}

\subsection*{Your turn}
\begin{itemize}
\raggedright
  \item \emph{Hear:} \emph{kalās}; YOU SAY: \textbf{\pa{ਕਲਾਸ}}
  \item \emph{Read it:} \textbf{\pa{ਕਲਾਸ}}
  \item \emph{Write it:} copy \textbf{\pa{ਕਲਾਸ}} once, checking the whole word
\end{itemize}

\begin{tcolorbox}[breakable,colback=teal!4,colframe=teal!35!black,title={Before you move on}]
Produce the school noun for “class” or “grade” without looking.
\end{tcolorbox}
\section[\texorpdfstring{\pa{ਕਲਾਸ} (kalaas)}{ਕਲਾਸ (kalaas)}]{Retrieve the new school word on its own}

\label{lesson:PA-C201-class-review}



\begin{quote}
\emph{Hear:} \emph{kalās}; YOU SAY: \textbf{\pa{ਕਲਾਸ}}
\end{quote}

\subsection*{Your turn}
\begin{itemize}
\raggedright
  \item \emph{Hear:} \emph{school class or grade}; YOU SAY: \textbf{\pa{ਕਲਾਸ}}
  \item \emph{Read it:} \textbf{\pa{ਕਲਾਸ}}
  \item \emph{Write it:} cover and write \textbf{\pa{ਕਲਾਸ}}
\end{itemize}

\begin{grammarlens}[title={Grammar Lens}]
Keep one school noun and the bounded class-or-grade meaning here. Wider forms, combinations, and other meanings stay outside this lesson.
\end{grammarlens}

\begin{tcolorbox}[breakable,colback=teal!4,colframe=teal!35!black,title={Before you move on}]
Say and then write the word once from memory.
\end{tcolorbox}
\section[\texorpdfstring{\pa{ਇਹ} \pa{ਕਲਾਸ} \pa{ਹੈ} (eh kalaas hai)}{ਇਹ ਕਲਾਸ ਹੈ (eh kalaas hai)}]{Put the new school noun into one identity frame}

\label{lesson:PA-C201-class-label}



\begin{quote}
Retrieve \textbf{\pa{ਇਹ}}, then the new school noun, then \textbf{\pa{ਹੈ}}.
\end{quote}

\begin{grammarlens}[title={Grammar Lens}]
\begin{quote}
\textbf{\pa{ਇਹ} \pa{ਕਲਾਸ} \pa{ਹੈ}} — \emph{eh kalās hai} — this is a class
\end{quote}

Keep the new whole noun between familiar \textbf{\pa{ਇਹ}} and \textbf{\pa{ਹੈ}}. Familiar \textbf{\pa{ਹੈ}} stays last in the already-owned identity frame. This chapter teaches only this one school label; it does not add a form rule or a wider construction.
\end{grammarlens}

\subsection*{Your turn}
\begin{itemize}
\raggedright
  \item \emph{Hear:} \emph{eh kalās hai}; YOU SAY: \textbf{\pa{ਇਹ} \pa{ਕਲਾਸ} \pa{ਹੈ}}
  \item \emph{Read it:} \textbf{\pa{ਇਹ} \pa{ਕਲਾਸ} \pa{ਹੈ}}
  \item \emph{Write it:} cover and write \textbf{\pa{ਇਹ} \pa{ਕਲਾਸ} \pa{ਹੈ}}
\end{itemize}

\begin{tcolorbox}[breakable,colback=teal!4,colframe=teal!35!black,title={Before you move on}]
Produce the complete identity frame.
\end{tcolorbox}
\section[\texorpdfstring{\pa{ਇਹ} \pa{ਕਲਾਸ} \pa{ਹੈ} (eh kalaas hai)}{ਇਹ ਕਲਾਸ ਹੈ (eh kalaas hai)}]{A bounded school-class payoff}

\label{lesson:PA-C201-payoff}



\begin{quote}
Retrieve the new school noun alone before rebuilding its identity frame.
\end{quote}

\subsection*{Your turn: four skills}
\begin{itemize}
\raggedright
  \item \emph{Hear:} \emph{eh kalās hai}; YOU SAY: \textbf{\pa{ਇਹ} \pa{ਕਲਾਸ} \pa{ਹੈ}}
  \item \emph{Read it:} \textbf{\pa{ਇਹ} \pa{ਕਲਾਸ} \pa{ਹੈ}}
  \item \emph{Write it:} cover the line and write \textbf{\pa{ਇਹ} \pa{ਕਲਾਸ} \pa{ਹੈ}}
\end{itemize}

\begin{tcolorbox}[breakable,colback=teal!4,colframe=teal!35!black,title={Before you move on}]
Say the new noun, then produce the complete identity frame.

This is the twenty-fourth prerequisite tranche exposed by Punjabi A2 mock drafting. It does not unblock Punjabi A2 mock 1 by itself. It is not an A2 mock, not A2 readiness or attainment, and not external or human validation.
\end{tcolorbox}
